\chapter{缓存、TLB 与页级性能}

CPU 核心很快，主内存相对很慢。缓存层级的存在，是为了让核心大多数时候从更近、更小、更快的存储中取数据。虚拟内存的存在，是为了让进程拥有独立地址空间、权限保护和按需映射。第一册已经讲过进程地址空间、映射和缺页的基本概念；本章不重复这些基础，而是继续追问：地址翻译、缓存映射、写共享、page walk、首次触页和 NUMA 放置怎样限制多核程序的吞吐。

\section{缓存层级与 cache line}

缓存不是按单个字节移动数据，而是按 cache line 移动。常见 cache line 大小是 64 字节。顺序访问数组时，一次 miss 可以带来后续多个元素；随机访问链表时，每个节点可能引入一次新的 miss。这个差异解释了为什么连续数组常常比指针结构更快。

缓存有容量、相联度和替换策略。容量不足会产生容量 miss，相联度不足会产生冲突 miss，多核心共享写会产生一致性流量。一个算法理论复杂度更低，不代表真实更快；如果它把连续访问变成随机访问，可能输给复杂度略高但局部性更好的算法。

从第二册的角度看，cache line 还是多核协作的最小一致性单元。两个线程哪怕写的是两个不同字段，只要字段落在同一条 cache line 上，硬件就必须在核心之间移动同一个一致性单元。很多并发程序的扩展性问题，不是算法复杂度问题，而是写入布局问题。后面讲锁、原子和无锁队列时，都要回到这个事实。

\begin{lstlisting}[numbers=none]
one cache line, 64 bytes in a common x86 machine

| counter[0] | counter[1] | counter[2] | counter[3] | ...
     T0           T1           T2           T3

looks independent in source code
shares one coherence unit in hardware
\end{lstlisting}

这类图要让读者形成一个直觉：源码里的变量边界，不等于硬件里的传输边界。硬件搬的是 cache line，TLB 记的是页，操作系统调度的是线程，分布式系统发送的是消息。边界错位，是系统性能问题最常见的来源之一。

\section{缓存映射和冲突}

缓存通常按 set associative 组织。一个地址会映射到某个 set，同一个 set 里只有有限个 way。如果多个热点地址恰好映射到同一个 set，而数量超过相联度，就会频繁互相驱逐，即使总工作集小于缓存容量也会慢。这类冲突 miss 很隐蔽，因为从容量角度看“不应该放不下”。

矩阵步长、数组对齐和结构体大小都可能触发冲突。若多个数组按相同大步长访问，并且基地址关系不佳，它们可能持续争用同一组 cache set。改变 padding、改变块大小或调整访问顺序，有时能显著改善性能。

\section{Inclusive、Exclusive 和共享缓存}

不同处理器的缓存层级策略不同。有的层级倾向 inclusive，也就是上层 cache 的内容也在下层保留；有的倾向 exclusive，尽量让不同层级存放不同内容；还有很多混合策略。软件通常不需要依赖某个具体策略，但要知道共享 LLC 不是无限资源。多个线程在同一 socket 上跑不同任务，可能互相驱逐 LLC。

这对并行任务调度有直接影响。把共享数据的任务放在共享缓存附近，可能复用 LLC；把互相干扰的大流式任务放在同一 socket，可能争抢带宽和 LLC。操作系统调度器会做一般性决策，但高性能程序常需要自己记录拓扑并做亲和性实验。

\section{写路径和 false sharing}

读缓存可以共享，写缓存必须取得 cache line 的独占所有权。两个线程写不同变量，如果这些变量落在同一条 cache line 上，硬件仍然会把它们当成同一块一致性单元反复迁移，这就是 false sharing。它不是锁竞争，却会表现出类似的扩展性下降。

解决 false sharing 的方法包括按线程分配独立槽位、使用对齐和填充、批量合并结果、减少共享写频率。并行归约通常先让每个线程写自己的 partial sum，最后再合并，正是为了避免每次加法都写同一个共享变量。

\begin{lstlisting}[language=C++,caption={用对齐槽位减少 false sharing}]
struct alignas(64) LocalCounter {
    std::int64_t value = 0;
};

std::vector<LocalCounter> counters(worker_count);
for (std::int32_t worker = 0; worker < worker_count; ++worker) {
    counters[static_cast<std::size_t>(worker)].value += 1;
}
\end{lstlisting}

这段代码不是说所有结构体都要强行 64 字节对齐，而是表达一个设计原则：如果多个线程会高频写不同槽位，这些槽位不要共享同一条 cache line。否则程序表面没有共享变量，硬件层面仍然在共享一致性单元。

false sharing 的危险在于它不像数据竞争那样直接破坏结果。程序可能完全正确，测试也全部通过，只是在核心数增加时吞吐突然停止增长。排查时要看写入频率、字段布局、线程到槽位的映射和 cache line 迁移事件。若没有 \cmd{perf c2c} 或类似工具，也可以用对齐填充做对照实验：只改变布局，不改变算法。如果对齐版本显著更快，就说明共享写路径值得重点检查。

\section{读共享和写共享}

共享数据不总是坏事。只读共享表、常量配置、不可变索引和代码段可以被多个核心同时缓存，通常扩展性很好。真正昂贵的是写共享，尤其是多个核心高频写同一条 cache line。读多写少的数据可以使用读写锁、RCU、版本化快照或 copy-on-write，目的都是让读路径尽量不进入写共享热路径。

写共享也分层次。最差情况是每次操作都写同一个全局变量，例如所有线程对一个原子计数器做 \code{fetch_add}。较好的情况是每个线程写本地槽位，周期性合并。更好的情况是让写入天然按数据分片分散，例如按 key 分片的哈希表或按 NUMA 节点分组的内存池。布局、同步和算法在这里是同一个问题的三个面。

\section{TLB 和页表}

虚拟地址访问物理内存前需要地址翻译。TLB 缓存虚拟页到物理页的映射。工作集页数太多或访问模式太随机时，TLB miss 会增加，核心需要走页表。大页可以减少页数，提高 TLB 覆盖范围，但也会影响内存分配、NUMA 放置和碎片。

理解 TLB 对数据结构设计很重要。一个跨很多小对象的图遍历，不只会造成 cache miss，还可能造成 TLB miss。紧凑存储、池化分配和按块遍历，可以同时改善 cache 和 TLB 行为。

TLB 的特殊之处在于，它限制的是“能高效覆盖多少虚拟页”。一个程序的工作集如果只看字节数不大，但分散在大量页面上，仍然可能受到 TLB miss 影响。链表、树、哈希表桶、对象池和图邻接结构都可能出现这个问题。把对象压缩到连续数组里，不只是减少 cache miss，也是在减少活跃页数。

\section{Page walk 的成本}

TLB miss 后，硬件需要按页表层级查找映射。x86-64 常见四级或五级页表，每一级页表项本身也在内存里，虽然页表项可能被缓存，但最坏情况下一个地址翻译会触发多次内存访问。若程序随机访问大量页面，真正拖慢它的不只是数据 load，还有地址翻译。

大页把页大小从 4 KiB 扩大到 2 MiB 或更大，显著增加 TLB 覆盖范围。它适合大数组、数据库缓存和数值计算工作集，但不适合所有场景。大页可能增加内存碎片，NUMA 放置也要更谨慎。使用大页前应先用计数器确认 TLB 确实是瓶颈。

Page walk 还会和缓存层级交织。页表项本身也要从内存层级读取，也可能命中或错过缓存。一个随机访问大数组的程序，表面上每次只读一个元素，实际上可能同时触发数据 cache miss 和页表访问。若工作集跨越大量页面，TLB miss 的成本会叠加在数据访问成本上。优化时不能只看“每个元素读几个字节”，还要看“每个时间窗口内活跃多少页”。

\section{缺页和内存分配}

分配内存不等于物理页已经准备好。操作系统通常按需分配页面，第一次触碰页面时可能发生 page fault。benchmark 如果把第一次触碰计入核心循环，会把页错误成本混入算法时间。严谨实验通常要预热、固定输入规模，并记录 page-faults。

缺页也分 minor fault 和 major fault。minor fault 不需要从磁盘读数据，可能只是建立页表映射；major fault 需要从存储设备读取，成本高得多。高性能计算通常要避免在热路径出现 major fault，也要避免把大量 minor fault 误判为算法本身慢。

第一册已经要求 benchmark 避免把初始化混进热路径。第二册要进一步要求初始化方式和计算方式一致。NUMA 机器上，如果主线程单独触碰全部页面，再让多个 socket 上的 worker 计算，页面可能集中在主线程所在节点，后续计算变成远端内存访问。正确实验通常让每个 worker 初始化自己之后要处理的分块，这样 page fault、物理页分配和计算拓扑保持一致。

\section{从虚拟内存到 mmap}

\code{mmap} 可以把文件或匿名内存映射到进程地址空间。它让文件 I/O 看起来像内存访问，但不意味着没有 I/O 成本。第一次访问未驻留页面时仍然会缺页，顺序扫描大文件时内核预读可能有效，随机访问时可能产生大量 page fault。

数据库、搜索引擎和日志系统常利用 \code{mmap} 或 page cache，但必须理解内核何时回收页面、何时写回脏页、何时产生 I/O 抖动。系统工程里，虚拟内存既是便利抽象，也是性能变量。

\section{Linux 源码阅读入口}

第一册的 Linux 源码阅读只要求读者能从用户态现象找到入口。第二册要读出数据结构和性能路径。围绕本章，可以固定一个 Linux 版本，沿下面几条路径做窄范围阅读：页错误入口可从体系结构相关的 fault handler 进入，再走到通用内存管理逻辑；VMA 管理可关注 \filepath{mm/mmap.c} 和相关 maple tree 结构；页表和 fault 处理可关注 \filepath{mm/memory.c}；NUMA 策略可关注 \filepath{mm/mempolicy.c}；page cache 和文件映射可从 \filepath{mm/filemap.c} 开始。

源码阅读报告不要写成“Linux 使用某某机制”这么宽泛。合格报告要写清内核版本、阅读路径、关键结构、用户态实验和结论边界。例如研究首次触页，可以写一个大数组实验，记录 page-faults，再阅读 fault 路径中如何建立匿名页映射。若当前设备权限不足，仍然可以保留源码阅读，但必须承认它不是当前运行内核的直接证据。

\section{本章的设计原则}

缓存和 TLB 优化可以压缩成四条原则。第一，让热路径连续访问，减少随机指针跳转。第二，让共享写分散，避免多个线程争用同一条 cache line。第三，让工作集按块复用，避免超出缓存后仍反复访问。第四，让页面放置和线程执行位置一致，避免首次触页造成远端访问。这四条原则会贯穿后面的数据布局、并行归约、线程池、工作窃取和分布式 shuffle。

\begin{warningbox}
不要把 \code{std::vector} 的连续性理解成所有成本都消失了。连续布局改善 cache line 和预取，但如果容量超过缓存，仍然会受内存带宽限制；如果首次访问大量新页，仍然会看到缺页成本。
\end{warningbox}

\begin{labbox}
实验：比较行优先访问矩阵和列优先访问矩阵。记录 cache-misses、dTLB 相关事件和时间。再把矩阵规模从 L1、L2、L3 可容纳范围逐步扩大，观察拐点。
\end{labbox}
